% ════════════════════════════════════════════════════════════════════════════
% GUÍA — LOGARITMOS II: PROPIEDADES Y CAMBIO DE BASE
% Semana 17 · 2026-S1 · Máximo + Martina · 2° Medio
% ════════════════════════════════════════════════════════════════════════════

\documentclass[11pt,a4paper]{article}
\usepackage[utf8]{inputenc}
\usepackage[T1]{fontenc}
\usepackage[spanish]{babel}
\usepackage{geometry}
\usepackage{amsmath,amssymb}
\usepackage{xcolor}
\usepackage{tcolorbox}
\usepackage{enumitem}
\usepackage{fancyhdr}
\usepackage{titlesec}
\usepackage{multirow}
\usepackage{array}
\usepackage{booktabs}
\usepackage{tikz}
\usepackage{hyperref}

\geometry{margin=2cm, top=2.5cm, bottom=2.5cm}

% ─── VARIABLES ───
\newcommand{\guiaTitulo}{Logaritmos II}
\newcommand{\guiaSubtitulo}{Propiedades y cambio de base}
\newcommand{\guiaUnidad}{Unidad de Álgebra}
\newcommand{\guiaCurso}{2° Medio --- Matemáticas}
\newcommand{\guiaAlumno}{\underline{\hspace{6cm}}}
\newcommand{\guiaSemana}{Semana 17}
\newcommand{\guiaSemestre}{2026-S1}

% ─── Colores ───
\definecolor{azulprincipal}{HTML}{2980B9}
\definecolor{azuloscuro}{HTML}{1A5276}
\definecolor{verdequimica}{HTML}{1ABC9C}
\definecolor{verdeclaro}{HTML}{D5F5E3}
\definecolor{naranja}{HTML}{E67E22}
\definecolor{rojo}{HTML}{E74C3C}
\definecolor{amarillorecuerda}{HTML}{FFF3CD}
\definecolor{amarilloborde}{HTML}{FFC107}
\definecolor{grisfondo}{HTML}{F5F5F5}

% ─── Cajas ───
\tcbset{
  recuerda/.style={
    colback=amarillorecuerda, colframe=amarilloborde,
    fonttitle=\bfseries, title={\color{black} RECUERDA},
    boxrule=1pt, arc=2mm, left=5pt, right=5pt
  },
  propiedad/.style={
    colback=verdeclaro!50, colframe=verdequimica!80!black,
    fonttitle=\bfseries, boxrule=0.8pt, arc=2mm, left=5pt, right=5pt
  },
  formula/.style={
    colback=grisfondo, colframe=azulprincipal,
    boxrule=0.8pt, arc=2mm, left=5pt, right=5pt
  },
  definicion/.style={
    colback=blue!5, colframe=azulprincipal,
    fonttitle=\bfseries, boxrule=0.8pt, arc=2mm, left=5pt, right=5pt
  },
  ejemplo/.style={
    colback=orange!5, colframe=naranja,
    fonttitle=\bfseries, boxrule=0.8pt, arc=2mm, left=5pt, right=5pt
  }
}

% ─── Encabezados ───
\pagestyle{fancy}
\fancyhf{}
\fancyhead[C]{\small\color{gray} \guiaTitulo{} --- \guiaCurso{}}
\fancyfoot[C]{\small\color{gray} Página \thepage}
\renewcommand{\headrulewidth}{0.4pt}

% ─── Títulos ───
\newcommand{\tituloseccion}[1]{%
  \noindent\colorbox{azulprincipal}{%
    \makebox[\dimexpr\textwidth-2\fboxsep\relax][l]{%
      \color{white}\normalfont\Large\bfseries\ SECCIÓN \thesection: #1%
    }%
  }%
}
\titleformat{\section}
  {}{}{0em}
  {\tituloseccion}
\titleformat{\subsection}
  {\normalfont\large\bfseries\color{azulprincipal}}
  {\thesubsection}{0.5em}{}

% ─── Niveles ───
\newcommand{\nivelbasico}{%
  \par\smallskip\colorbox{green!60!black}{\color{white}\small\bfseries\ NIVEL BÁSICO\ }\par\smallskip}
\newcommand{\nivelintermedio}{%
  \par\smallskip\colorbox{naranja}{\color{white}\small\bfseries\ NIVEL INTERMEDIO\ }\par\smallskip}
\newcommand{\nivelavanzado}{%
  \par\smallskip\colorbox{rojo}{\color{white}\small\bfseries\ NIVEL AVANZADO\ }\par\smallskip}

% ─── Líneas ───
\newcommand{\lineasrespuesta}[1]{%
  \par\vspace{2mm}%
  \foreach \i in {1,...,#1}{\noindent\rule{\textwidth}{0.3pt}\par\vspace{5mm}}%
  \vspace{2mm}%
}

\begin{document}

% ════════════════════════
% PORTADA
% ════════════════════════
\thispagestyle{empty}
\vspace*{3cm}
\begin{center}
  {\Huge\bfseries\color{azulprincipal} GUÍA DE TRABAJO}\\[8pt]
  {\LARGE\bfseries\color{azuloscuro} \guiaTitulo}\\[15pt]
  {\color{azulprincipal}\rule{8cm}{2pt}}\\[15pt]
  {\Large \guiaUnidad{} \textbar{} \guiaCurso}\\[8pt]
  {\large Alumno/a: \guiaAlumno}\\[4pt]
  {\large \guiaSemana{} --- \guiaSemestre}\\[25pt]
  {\itshape\color{gray}
  \guiaSubtitulo}
\end{center}

\vspace{2cm}
\noindent\textbf{Contenidos:}
\begin{enumerate}[leftmargin=2cm]
  \item Repaso flash: definición y casos especiales de la semana 16.
  \item Propiedad del producto (logaritmo de un producto).
  \item Propiedad del cociente (logaritmo de un cociente).
  \item Propiedad de la potencia (logaritmo de una potencia).
  \item Cambio de base.
  \item Ejercicios por nivel (básico, intermedio, avanzado).
\end{enumerate}
\newpage

% ════════════════════════
% SECCIÓN 1
% ════════════════════════
\section{REPASO FLASH DE LA SEMANA 16}

\begin{tcolorbox}[recuerda]
\textbf{Definición:} $\log_b(a) = c \;\iff\; b^{c} = a$, con $b > 0$, $b \ne 1$, $a > 0$.

\medskip
\textbf{Casos especiales:}
\begin{itemize}[nosep]
  \item $\log_b(1) = 0$
  \item $\log_b(b) = 1$
  \item $\log_b(b^{n}) = n$
\end{itemize}
\end{tcolorbox}

\textbf{Calentamiento} (resolver mentalmente):
\begin{enumerate}[label=\alph*),leftmargin=1cm,itemsep=4pt]
  \item $\log_2(32) = $ \rule{2cm}{0.4pt}
  \item $\log_3(1) = $ \rule{2cm}{0.4pt}
  \item $\log_5(5^{7}) = $ \rule{2cm}{0.4pt}
  \item $\log_{10}(10\,000) = $ \rule{2cm}{0.4pt}
  \item $\log_4(4) = $ \rule{2cm}{0.4pt}
\end{enumerate}

{\small\textit{Respuestas: a) $5$ \; b) $0$ \; c) $7$ \; d) $4$ \; e) $1$.}}

\newpage

% ════════════════════════
% SECCIÓN 2
% ════════════════════════
\section{LAS TRES PROPIEDADES OPERATORIAS}

Las tres propiedades del logaritmo son \textbf{consecuencia directa} de las propiedades de las potencias. La clave es esta correspondencia:

\begin{tcolorbox}[formula]
\begin{center}
\begin{tabular}{@{}ccc@{}}
\textbf{Propiedad de potencias} & & \textbf{Propiedad de logaritmos} \\[4pt]
$b^{m} \cdot b^{n} = b^{m+n}$ & $\longleftrightarrow$ & $\log_b(M \cdot N) = \log_b M + \log_b N$ \\[6pt]
$\dfrac{b^{m}}{b^{n}} = b^{m-n}$ & $\longleftrightarrow$ & $\log_b\!\left(\dfrac{M}{N}\right) = \log_b M - \log_b N$ \\[8pt]
$(b^{m})^{n} = b^{m \cdot n}$ & $\longleftrightarrow$ & $\log_b(M^{n}) = n \cdot \log_b M$ \\
\end{tabular}
\end{center}
\end{tcolorbox}

\subsection{Propiedad 1: Logaritmo de un producto}

\begin{tcolorbox}[propiedad, title={Propiedad del producto}]
\[
  \boxed{\;\log_b(M \cdot N) = \log_b M + \log_b N\;}
\]
``El logaritmo de un \textbf{producto} es la \textbf{suma} de los logaritmos.''
\end{tcolorbox}

\begin{tcolorbox}[definicion, title={¿De dónde sale?}]
Sean $\log_b M = p$ y $\log_b N = q$. Entonces $M = b^{p}$ y $N = b^{q}$.
\[
  M \cdot N = b^{p} \cdot b^{q} = b^{p+q}
\]
Tomando logaritmo en base $b$ de ambos lados:
\[
  \log_b(M \cdot N) = \log_b(b^{p+q}) = p + q = \log_b M + \log_b N \qquad \checkmark
\]
\end{tcolorbox}

\begin{tcolorbox}[ejemplo, title={Ejemplo: $\log_2(8 \cdot 16)$}]
\textbf{Método directo:} $8 \cdot 16 = 128 = 2^{7}$, así que $\log_2(128) = 7$.

\textbf{Con la propiedad:} $\log_2(8 \cdot 16) = \log_2 8 + \log_2 16 = 3 + 4 = 7$. \checkmark
\end{tcolorbox}

\subsection{Propiedad 2: Logaritmo de un cociente}

\begin{tcolorbox}[propiedad, title={Propiedad del cociente}]
\[
  \boxed{\;\log_b\!\left(\frac{M}{N}\right) = \log_b M - \log_b N\;}
\]
``El logaritmo de un \textbf{cociente} es la \textbf{resta} de los logaritmos.''
\end{tcolorbox}

\begin{tcolorbox}[definicion, title={¿De dónde sale?}]
Con $M = b^{p}$ y $N = b^{q}$:
\[
  \frac{M}{N} = \frac{b^{p}}{b^{q}} = b^{p-q} \;\;\Longrightarrow\;\; \log_b\!\left(\frac{M}{N}\right) = p - q = \log_b M - \log_b N \qquad \checkmark
\]
\end{tcolorbox}

\begin{tcolorbox}[ejemplo, title={Ejemplo: $\log_3\!\left(\frac{81}{9}\right)$}]
\textbf{Directo:} $81 \div 9 = 9 = 3^{2}$, así que $\log_3(9) = 2$.

\textbf{Con la propiedad:} $\log_3 81 - \log_3 9 = 4 - 2 = 2$. \checkmark
\end{tcolorbox}

\subsection{Propiedad 3: Logaritmo de una potencia}

\begin{tcolorbox}[propiedad, title={Propiedad de la potencia}]
\[
  \boxed{\;\log_b(M^{n}) = n \cdot \log_b M\;}
\]
``El \textbf{exponente} del argumento \textbf{baja y multiplica}.''
\end{tcolorbox}

\begin{tcolorbox}[definicion, title={¿De dónde sale?}]
Con $M = b^{p}$:
\[
  M^{n} = (b^{p})^{n} = b^{pn} \;\;\Longrightarrow\;\; \log_b(M^{n}) = pn = n \cdot \log_b M \qquad \checkmark
\]
\end{tcolorbox}

\begin{tcolorbox}[ejemplo, title={Ejemplo: $\log_2(8^{4})$}]
\textbf{Con la propiedad:} $\log_2(8^{4}) = 4 \cdot \log_2 8 = 4 \cdot 3 = 12$.

\textbf{Verificación:} $8^{4} = (2^{3})^{4} = 2^{12}$, y $\log_2(2^{12}) = 12$. \checkmark
\end{tcolorbox}

\begin{tcolorbox}[recuerda]
\textbf{Errores frecuentes:}
\begin{itemize}[nosep]
  \item $\log_b(M + N) \;\ne\; \log_b M + \log_b N$. \quad La propiedad es para \textbf{productos}, no sumas.
  \item $\log_b(M \cdot N) \;\ne\; \log_b M \cdot \log_b N$. \quad Producto $\to$ \textbf{suma}, no producto.
  \item $(\log_b M)^{n} \;\ne\; n \cdot \log_b M$. \quad El exponente debe estar \textbf{en el argumento}, no afuera.
\end{itemize}
\end{tcolorbox}

\newpage

% ════════════════════════
% SECCIÓN 3
% ════════════════════════
\section{CAMBIO DE BASE}

¿Cuánto vale $\log_2(10)$? No es un entero (porque $2^3 = 8$ y $2^4 = 16$). La calculadora solo tiene $\log_{10}$ (tecla \texttt{log}) y $\log_e$ (tecla \texttt{ln}). Necesitamos una forma de calcular un logaritmo en cualquier base usando una base que sí tengamos.

\begin{tcolorbox}[propiedad, title={Fórmula de cambio de base}]
\[
  \boxed{\;\log_b(a) = \frac{\log_c(a)}{\log_c(b)}\;}
\]
Funciona para \textbf{cualquier} base $c > 0$, $c \ne 1$. En la práctica, usamos $c = 10$ o $c = e$.
\end{tcolorbox}

\begin{tcolorbox}[definicion, title={¿De dónde sale?}]
Sea $\log_b(a) = x$. Entonces $b^{x} = a$.

Tomo $\log_c$ de ambos lados:
\[
  \log_c(b^{x}) = \log_c(a) \;\;\Longrightarrow\;\; x \cdot \log_c(b) = \log_c(a) \;\;\Longrightarrow\;\; x = \frac{\log_c(a)}{\log_c(b)} \qquad \checkmark
\]
\end{tcolorbox}

\begin{tcolorbox}[ejemplo, title={Ejemplo: $\log_2(10)$ con calculadora}]
\[
  \log_2(10) = \frac{\log_{10}(10)}{\log_{10}(2)} = \frac{1}{0{,}3010\ldots} \approx 3{,}322
\]
\textbf{Verificación:} $2^{3{,}322} \approx 10$. \checkmark

Otra forma equivalente con $\ln$:
\[
  \log_2(10) = \frac{\ln(10)}{\ln(2)} = \frac{2{,}3026\ldots}{0{,}6931\ldots} \approx 3{,}322
\]
\end{tcolorbox}

\begin{tcolorbox}[ejemplo, title={Ejemplo: $\log_5(80)$}]
\[
  \log_5(80) = \frac{\log(80)}{\log(5)} = \frac{1{,}9031\ldots}{0{,}6990\ldots} \approx 2{,}723
\]
Verificación: $5^{2} = 25$ y $5^{3} = 125$, así que un valor entre $2$ y $3$ tiene sentido. \checkmark
\end{tcolorbox}

\newpage

% ════════════════════════
% SECCIÓN 4 — EJERCICIOS
% ════════════════════════
\section{EJERCICIOS}

\nivelbasico

\textbf{(1)} Expande usando las propiedades (no calcules el valor numérico).
\begin{enumerate}[label=\alph*),leftmargin=1cm,itemsep=6pt]
  \item $\log_3(9 \cdot 27) = $ \rule{5cm}{0.4pt}
  \item $\log_2\!\left(\dfrac{32}{4}\right) = $ \rule{5cm}{0.4pt}
  \item $\log_5(25^{3}) = $ \rule{5cm}{0.4pt}
  \item $\log_{10}(100 \cdot 1000) = $ \rule{5cm}{0.4pt}
\end{enumerate}

\medskip
\textbf{(2)} Condensa en un solo logaritmo.
\begin{enumerate}[label=\alph*),leftmargin=1cm,itemsep=6pt]
  \item $\log_2 5 + \log_2 3 = $ \rule{5cm}{0.4pt}
  \item $\log_3 12 - \log_3 4 = $ \rule{5cm}{0.4pt}
  \item $3 \cdot \log_2 5 = $ \rule{5cm}{0.4pt}
\end{enumerate}

\medskip
\textbf{(3)} Calcula el valor numérico final (expandir con propiedades + usar la definición).
\begin{enumerate}[label=\alph*),leftmargin=1cm,itemsep=6pt]
  \item $\log_2(4 \cdot 8)$
  \lineasrespuesta{2}
  \item $\log_3\!\left(\dfrac{81}{3}\right)$
  \lineasrespuesta{2}
  \item $\log_2(16^{2})$
  \lineasrespuesta{2}
\end{enumerate}

\nivelintermedio

\textbf{(4)} Expande completamente.
\begin{enumerate}[label=\alph*),leftmargin=1cm,itemsep=6pt]
  \item $\log_2(8x^{3})$ \hfill \textit{Pista: producto primero, potencia después.}
  \lineasrespuesta{2}
  \item $\log_b\!\left(\dfrac{a^{2}}{c}\right)$
  \lineasrespuesta{2}
  \item $\log_3\!\left(\dfrac{9x^{4}}{y^{2}}\right)$
  \lineasrespuesta{3}
\end{enumerate}

\medskip
\textbf{(5)} Condensa en un solo logaritmo.
\begin{enumerate}[label=\alph*),leftmargin=1cm,itemsep=6pt]
  \item $2\log_b x + 3\log_b y$
  \lineasrespuesta{2}
  \item $\log_5 a - \tfrac{1}{2}\log_5 b$
  \lineasrespuesta{2}
  \item $\log_2 3 + \log_2 5 - \log_2 15$
  \lineasrespuesta{2}
\end{enumerate}

\medskip
\textbf{(6)} Calcula con cambio de base (usa $\log_{10}$ o $\ln$; puedes usar calculadora).
\begin{enumerate}[label=\alph*),leftmargin=1cm,itemsep=6pt]
  \item $\log_3(20) \approx$ \rule{3cm}{0.4pt}
  \item $\log_8(50) \approx$ \rule{3cm}{0.4pt}
  \item $\log_2(100) \approx$ \rule{3cm}{0.4pt}
\end{enumerate}

\nivelavanzado

\textbf{(7)} Demuestra la propiedad del producto sin mirar la guía. Parte de ``sean $\log_b M = p$ y $\log_b N = q$'' y llega a $\log_b(MN) = p + q$.
\lineasrespuesta{5}

\textbf{(8)} Demuestra la fórmula de cambio de base: parte de $\log_b(a) = x$ y llega a $x = \dfrac{\log_c(a)}{\log_c(b)}$.
\lineasrespuesta{5}

\textbf{(9)} Simplifica sin calculadora.
\begin{enumerate}[label=\alph*),leftmargin=1cm,itemsep=6pt]
  \item $\log_2 6 - \log_2 3 + \log_2 8$
  \lineasrespuesta{3}
  \item $\dfrac{\log_9 27}{\log_9 3}$ \hfill \textit{Pista: ¿hay un cambio de base escondido?}
  \lineasrespuesta{3}
  \item $\log_4 8$ \hfill \textit{Pista: escribir todo como potencias de 2.}
  \lineasrespuesta{3}
\end{enumerate}

\newpage

% ════════════════════════
% RESPUESTAS
% ════════════════════════
\section*{Respuestas seleccionadas}
\addcontentsline{toc}{section}{Respuestas seleccionadas}

\begin{tcolorbox}[propiedad, title=Soluciones]
\textbf{(1)} a) $\log_3 9 + \log_3 27 = 2 + 3 = 5$; \; b) $\log_2 32 - \log_2 4 = 5 - 2 = 3$; \; c) $3 \cdot \log_5 25 = 3 \cdot 2 = 6$; \; d) $\log 100 + \log 1000 = 2 + 3 = 5$.

\medskip
\textbf{(2)} a) $\log_2(15)$; \; b) $\log_3(3) = 1$; \; c) $\log_2(125)$.

\medskip
\textbf{(3)} a) $\log_2 4 + \log_2 8 = 2 + 3 = 5$; \; b) $\log_3 81 - \log_3 3 = 4 - 1 = 3$; \; c) $2 \cdot \log_2 16 = 2 \cdot 4 = 8$.

\medskip
\textbf{(4)} a) $\log_2 8 + 3\log_2 x = 3 + 3\log_2 x$; \; b) $2\log_b a - \log_b c$; \; c) $2 + 4\log_3 x - 2\log_3 y$.

\medskip
\textbf{(5)} a) $\log_b(x^{2}y^{3})$; \; b) $\log_5\!\left(\dfrac{a}{\sqrt{b}}\right)$; \; c) $\log_2\!\left(\dfrac{3 \cdot 5}{15}\right) = \log_2 1 = 0$.

\medskip
\textbf{(6)} a) $\approx 2{,}727$; \; b) $\approx 1{,}881$; \; c) $\approx 6{,}644$.

\medskip
\textbf{(9)} a) $\log_2\!\left(\dfrac{6 \cdot 8}{3}\right) = \log_2 16 = 4$; \; b) $\log_3 27 = 3$; \; c) $\dfrac{3}{2}$.
\end{tcolorbox}

\begin{tcolorbox}[recuerda]
Los ejercicios 7 y 8 (demostraciones) se revisan en clase.
\end{tcolorbox}

\end{document}
